\section*{अध्याय \ref{ch:g10:earth-energy-budget} --- पृथ्वी का ऊर्जा बजट}

\begin{solution}{exo:g10:earth-energy-budget:1}
किसी स्थान का \omterm{def:g10:earth-energy-budget:insolation}{सूर्यातप} उसके ऊपर वायुमंडल के शीर्ष पर एक क्षैतिज वर्ग मीटर द्वारा प्राप्त
सूर्य-प्रकाश की शक्ति है (एक दिन या एक वर्ष पर औसत)। किसी सतह का \omterm{def:g10:earth-energy-budget:albedo}{एल्बिडो} उस पर पड़ने
वाले सूर्य-प्रकाश का वह भाग है जिसे वह परावर्तित करती है। ताज़ा हिम (0.8 से 0.9) खुले
महासागर (0.06) से कहीं अधिक परावर्तित करता है।
\end{solution}

\begin{solution}{exo:g10:earth-energy-budget:2}
परावर्तित: $100/340 \approx \qty{29}{\percent}$; हवा द्वारा अवशोषित:
$80/340 \approx \qty{24}{\percent}$; सतह द्वारा अवशोषित: $160/340
\approx \qty{47}{\percent}$।
\end{solution}

\begin{solution}{exo:g10:earth-energy-budget:3}
डामर: $800 \times (1 - 0.04) = \qty{768}{W}$। हिम: $800 \times (1 -
0.85) = \qty{120}{W}$, छह गुना से भी कम।
\end{solution}

\begin{solution}{exo:g10:earth-energy-budget:4}
वाष्पन: \qty{82}{W/m^2}, जबकि हवा के गरम होने के लिए \qty{21}{W/m^2}। वाष्पित होने वाला
जल अपनी \omterm{def:g10:earth-energy-budget:latent-heat}{गुप्त ऊष्मा} समुद्र से लेता है; वाष्प हवा के साथ ऊपर उठती है, और जहाँ वह बादल
की बूँदों में संघनित होती है, अक्सर किलोमीटरों ऊपर, वहाँ वह यह ऊर्जा आसपास की हवा को लौटा
देती है।
\end{solution}

\begin{solution}{exo:g10:earth-energy-budget:5}
दिसंबर में उत्तरी गोलार्ध सूर्य से दूर झुका होता है: ध्रुव पर सूर्य पूरे दिन क्षितिज के नीचे
रहता है, इसलिए कोई सूर्य-प्रकाश नहीं पहुँचता। दिसंबर अयनांत पर
$90^\circ - 23.44^\circ =
\qty{66.56}{\degree}$ उत्तर से ध्रुव की ओर हर जगह ऐसा होता है।
\end{solution}

\begin{solution}{exo:g10:earth-energy-budget:6}
$1361 \sin 90^\circ = \qty{1361}{W}$; $1361 \sin 30^\circ =
\qty{680.5}{W}$; $1361 \sin 10^\circ \approx \qty{236}{W}$। $S_0$ के आधे के लिए
$\sin h = 0.5$ चाहिए, यानी $h = \qty{30}{\degree}$।
\end{solution}

\begin{solution}{exo:g10:earth-energy-budget:7}
परावर्तित 100, हवा द्वारा अवशोषित 80, सतह द्वारा अवशोषित 160:
$100 + 80 + 160 = 340$, सारे सूर्य-प्रकाश का हिसाब मिल जाता है। सतह पर,
$185 = 25 + 160$: जो परावर्तित नहीं होता वह अवशोषित होता है।
\end{solution}

\begin{solution}{exo:g10:earth-energy-budget:8}
चकती: $\pi \times (6.371\times10^6)^2 \approx \qty{1.275e14}{m^2}$; शक्ति:
$1361 \times 1.275\times10^{14} \approx \qty{1.74e17}{W}$। गोला: $4\pi
R^2 \approx \qty{5.10e14}{m^2}$। औसत: $1.74\times10^{17}/5.10\times
10^{14} \approx \qty{340}{W/m^2}$, $S_0$ का एक-चौथाई।
\end{solution}

\begin{solution}{exo:g10:earth-energy-budget:9}
$0.7 \times 5.10\times10^{14} \times 3.16\times10^7 \approx
\qty{1.1e22}{J}$ एक वर्ष में।
\end{solution}

\begin{solution}{exo:g10:earth-energy-budget:10}
बादल $100 - 53 = \qty{47}{W/m^2}$ अधिक सूर्य-प्रकाश परावर्तित करते हैं, जो इसलिए
अवशोषित नहीं होता। वे $267 - 239 = \qty{28}{W/m^2}$ अवरक्त रोक लेते हैं। कुल मिलाकर
पृथ्वी $47 - 28 = \qty{19}{W/m^2}$ खोती है: बादल,
समग्र रूप से, उसे ठंडा करते हैं।
\end{solution}

\begin{solution}{exo:g10:earth-energy-budget:11}
लगभग 10 मई से जुलाई के अंत तक, यानी मई, जून और जुलाई में। भूमध्य
रेखा पर सूर्य दोनों विषुवों, मार्च और सितंबर, पर दोपहर को सिर के ऊपर से गुज़रता है;
अयनांतों पर वह शीर्षबिंदु से \qty{23.44}{\degree} दूर से गुज़रता है, इसलिए भूमध्य रेखा
जून और दिसंबर में थोड़ा कम पाती है।
\end{solution}

\begin{solution}{exo:g10:earth-energy-budget:12}
हर सेकंड: $82/(2.47\times10^6) \approx \qty{3.32e-5}{kg}$। हर दिन:
$3.32\times10^{-5} \times 86\,400 \approx \qty{2.9}{kg}$, लगभग
\qty{2.9}{mm} की परत। हर वर्ष: $3.32\times10^{-5} \times 3.16\times10^7 \approx
\qty{1050}{kg}$, लगभग \qty{1.05}{m} जल।
\end{solution}

\begin{solution}{exo:g10:earth-energy-budget:13}
$S_0 \sin\delta = S_0/\pi$ से $\sin\delta = 1/\pi \approx 0.318$ मिलता है, इसलिए
$\delta \approx \qty{18.6}{\degree}$। तब $\sin\lambda = 0.318/\sin 23.44^\circ
\approx 0.318/0.398 \approx 0.800$: $\lambda \approx \qty{53.15}{\degree}$ या
$180^\circ - 53.15^\circ = \qty{126.85}{\degree}$। इन दोनों के बीच $\lambda$
\qty{73.7}{\degree} बढ़ता है, यानी वर्ष का $73.7/360$ भाग, लगभग
\num{75} दिन, मई के मध्य से जुलाई के अंत तक।
\end{solution}

\begin{solution}{exo:g10:earth-energy-budget:14}
बर्फ़ सूर्य-प्रकाश का $1 - 0.6 = 0.40$ भाग अवशोषित करती है, महासागर $1 - 0.06 =
0.94$। अतिरिक्त शक्ति: $(0.94 - 0.40) \times 300 \times 10^{12} \approx
\qty{1.6e14}{W}$, ग्रह के सारे प्रकाशसंश्लेषण से अधिक। अतिरिक्त
ऊर्जा जल को गरम करती है, जो हिम-खंडों के किनारों पर और बर्फ़ पिघलाता है,
जिससे और गहरे रंग का जल खुलता है: परिवर्तन स्वयं को पोषित करता है।
\end{solution}

\begin{solution}{exo:g10:earth-energy-budget:15}
सबसे निकट: $1361 \times (149.6/147.1)^2 \approx \qty{1408}{W/m^2}$; सबसे दूर:
$1361 \times (149.6/152.1)^2 \approx \qty{1317}{W/m^2}$; लगभग
\qty{7}{\percent} का अंतर। \qty{60}{\degree} उत्तर पर \omterm{def:g10:earth-energy-budget:insolation}{सूर्यातप} जून में दिसंबर से बीस
गुना अधिक है: कुछ प्रतिशत का परिवर्तन इसका कारण नहीं हो सकता। इसके अलावा पृथ्वी जनवरी
के आरंभ में, उत्तरी सर्दी के दौरान, सूर्य के सबसे निकट होती है, और दोनों गोलार्धों में, जो
सूर्य से एक ही दूरी पर हैं, उलटी ऋतुएँ होती हैं: ऋतुएँ अक्ष का झुकाव बनाता है, दूरी
नहीं।
\end{solution}

\begin{solution}{pb:g10:earth-energy-budget:1}
\textbf{1.} $\pi \times (6.371\times10^6)^2 \approx \qty{1.275e14}{m^2}$।

\textbf{2.} $1361 \times 1.275\times10^{14} \approx \qty{1.735e17}{W}$।

\textbf{3.} $4\pi R^2 \approx \qty{5.101e14}{m^2}$।

\textbf{4.} $1.735\times10^{17}/5.101\times10^{14} \approx \qty{340}{W/m^2}$,
यानी $S_0/4$।

\textbf{5.} भूमध्य रेखा: \qty{1361}{W}; \qty{60}{\degree} उत्तर: $1361 \sin
30^\circ = \qty{680.5}{W}$।

\textbf{6.} $100/340 \approx 0.294$।

\textbf{7.} $340 - 100 = \qty{240}{W/m^2}$; पूरी पृथ्वी पर, $240 \times
5.101\times10^{14} \approx \qty{1.224e17}{W}$।

\textbf{8.} $53/340 \approx 0.156$। बादल \omterm{def:g10:earth-energy-budget:albedo}{एल्बिडो} को
$100/53 \approx 1.9$ के गुणक से बढ़ाते हैं: लगभग दुगुना कर देते हैं।

\textbf{9.} $(240/5.67\times10^{-8})^{1/4} \approx \qty{255}{K}$, सतह के
\qty{288}{K} से लगभग \qty{33}{K} नीचे।

\textbf{10.} $160 + 342 = \qty{502}{W/m^2}$।

\textbf{11.} $82 + 21 + 398 = \qty{501}{W/m^2}$। सतह लगभग
\qty{1}{W/m^2} पाती है: पूर्णांकित \qty{0.7}{W/m^2} का असंतुलन।

\textbf{12.} $342/160 \approx 2.1$। हवा सतह को अवरक्त के रूप में उस सूर्य-प्रकाश के
दुगुने से भी अधिक ऊर्जा भेजती है जिसे सतह अवशोषित करती है: ग्रीनहाउस
प्रभाव।

\textbf{13.} $82 \times 3.156\times10^7/(2.47\times10^6) \approx
\qty{1050}{kg}$, लगभग \qty{1.05}{m} मोटी परत।

\textbf{14.} $0.7 \times 5.101\times10^{14} \approx \qty{3.6e14}{W}$।

\textbf{15.} $2805/72.07 \approx \qty{38.9}{kJ}$ प्रति ग्राम कार्बन।

\textbf{16.} $104.9\times10^9$ टन, यानी $1.049\times10^{17}$~g;
$1.049\times10^{17} \times 38.9\times10^3 \approx \qty{4.08e21}{J}$।

\textbf{17.} $4.08\times10^{21}/3.156\times10^7 \approx \qty{1.29e14}{W}$।

\textbf{18.} $1.29\times10^{14}/5.101\times10^{14} \approx
\qty{0.25}{W/m^2}$।

\textbf{19.} $0.7/0.25 \approx 2.8$: पृथ्वी का असंतुलन प्रकाशसंश्लेषण द्वारा संचित
शक्ति का लगभग तीन गुना है।

\textbf{20.} $1.29\times10^{14}/1.224\times10^{17} \approx 0.0011$:
प्रकाशसंश्लेषण पृथ्वी द्वारा अवशोषित सूर्य-प्रकाश का \textbf{लगभग \qty{0.1}{\percent}}
संचित करता है --- एक हज़ारवाँ भाग।
\end{solution}
